\section{Вывод}

Выполнив лабораторную работу, мы изучили принцип работы мультивибратора, триггера и кипп-реле и познакомились с основными явлениями, происходящими в системах с разрывными колебаниями.

\begin{enumerate}
	\item В режиме мультивибратора наблюдали осциллограммы напряжения и тока. Период колебаний тока совпадает с периодом колебаний напряжения, а амплитуды составляют:
	\[
	U_{\max} = 1000~\text{мВ}, \qquad U_{\min} = 650~\text{мВ}, \qquad T = 170~\text{мкс}.
	\]
	\item В режиме триггера длительность импульса на выходе совпадает с периодом повторения запускающих импульсов:
	\[
	T_\text{тр} = T_\text{и} = 1~\text{мс}.
	\]
	Схема перестаёт работать как триггер при слишком коротком (7--20~мкс) и слишком длинном (20--37~мкс) запускающем импульсе: минимальная допустимая длительность лежит между 7 и 20~мкс, максимальная --- между 20 и 37~мкс. Порог по амплитуде запускающего импульса составляет 0,5--0,9~В снизу и 0,9--1,5~В сверху. При частоте $f_\text{и} = 11{,}8$~кГц, амплитуде $A = 900$~мВ и заполнении 11~\% получено деление частоты: переключение происходит на каждый второй запускающий импульс.
	\item В режиме кипп-реле короткий запускающий импульс формирует на выходе широкий прямоугольный импульс, длительность которого определяется только параметрами схемы. Границы рабочего диапазона:
	\[
	\tau_{\text{и}\,\min} \approx 0{,}08~\text{мкс}, \qquad \tau_{\text{и}\,\max} \approx 2{,}8~\text{мкс}, \qquad A_{\min} = 0{,}21~\text{В}, \qquad A_{\max} \approx 0{,}99~\text{В}.
	\]
\end{enumerate}
